\hypertarget{appendix-f8657bac66e5}{}
\section{Matrix norms and numerical error}\label{appendix-f8657bac66e5-matrix-norms-and-numerical-error}

The spectral norm is the largest 2-norm stretch and equals the largest singular value. The Frobenius norm measures aggregate squared entry energy. Residuals, orthogonality errors, and condition numbers need a declared norm and scale; a small residual alone need not imply a sensitive solution is accurate.

\textbf{Use it when:} setting scaled tolerances or interpreting numerical diagnostics.

\textbf{Developed in Phase 8:} The spectral norm measures worst-case amplification, while the Frobenius norm accumulates squared energy across entries. The chapter uses them to scale residuals, reconstruction errors, and orthogonality defects, and separates backward error in a nearby problem from forward error in the requested answer. \href{../01-Foundations-I/main.pdf\#nameddest=01-foundations-i-ch019-matrix-norms-put-scale-into-error-statements}{norms and scale (01-Foundations-I, 01-foundations-i-ch019-matrix-norms-put-scale-into-error-statements)} · \href{../01-Foundations-I/main.pdf\#nameddest=01-foundations-i-ch020-residuals-distinguish-backward-accuracy-from-forward-certainty}{residuals and error interpretation (01-Foundations-I, 01-foundations-i-ch020-residuals-distinguish-backward-accuracy-from-forward-certainty)}

\subsection{Prerequisites}\label{appendix-f8657bac66e5-prerequisites}

\hyperref[appendix-49933872bb51]{Norms lengths and induced distances}, \hyperref[appendix-1fad5e7f36fa]{Singular value decomposition}, \href{../01-Foundations-I/main.pdf\#nameddest=appendix-a7fbaf6d1b1c}{Forward and backward error (01-Foundations-I, appendix-a7fbaf6d1b1c)}.

\subsection{Used by}\label{appendix-f8657bac66e5-used-by}

\hyperref[appendix-541306fdbc1b]{Moore-Penrose pseudoinverse}, \hyperref[appendix-0f4bc2bf2b1f]{Matrix functions}, \hyperref[appendix-67c87a3d9e94]{Sparse and iterative solvers}.

\textbf{Teaching path:} \href{../01-Foundations-I/main.pdf\#nameddest=01-foundations-i-ch019-matrix-norms-give-scale-to-residuals-and-error-claims}{norms and error (01-Foundations-I, 01-foundations-i-ch019-matrix-norms-give-scale-to-residuals-and-error-claims)}.

\subsection{A matrix norm can protect a geometry bound}\label{appendix-f8657bac66e5-a-matrix-norm-can-protect-a-geometry-bound}

A local sphere under a linear transform is generally an ellipsoid. The Frobenius norm times its original radius gives a safe, possibly loose enclosing radius. Taking only the longest matrix column can fail under shear. \href{../04-Visual-Computing-II/main.pdf\#nameddest=04-visual-computing-ii-ch028-a-line-corridor-can-be-a-cheap-prefilter-for-shadow-caster-work}{The culling example (04-Visual-Computing-II, 04-visual-computing-ii-ch028-a-line-corridor-can-be-a-cheap-prefilter-for-shadow-caster-work)} derives the bound, demonstrates the shear counterexample, and links its geometric figure. This connects numerical matrix norms to conservative rejection decisions.
